\textit{Jailbreaking} refers to disabling or circumventing restrictions placed upon a technology by its manufacturer or distributor, usually for the purpose of accessing the full range of the technology's capabilities. The term became popularized among technologist communities in the late 2000s and early 2010s and referred primarily to accessing and modifying the iPhone's filesystem to allow arbitrary applications and programs to run on the phone's hardware, bypassing copyright protection measures \cite{cheng_iphone_2010,wolk_iphone_2010}. We refer to this style of jailbreak as a \textit{classical jailbreak.} Classical jailbreaking is widely associated with hacker-culture ideals of free open source software (FOSS) and end-user freedom \cite{mohammed_what_2021, gewirtz_your_2026}; the jailbreak allows the user to access the ``true'' technology beneath the manufacturer's ``walled garden'' of approved use-cases \cite{mcsherry_breaking_2010}. 

The LLM jailbreaking community has adopted the term jailbreaking for a new paradigm of exploits, claiming similar cultural values and end-user aims. \textit{Persona-based} jailbreaks are a particularly effective class of jailbreaks \cite{shah_scalable_2023, zhang_enhancing_2025, yu_dont_2024} which create constructed identities designed to bypass model safety and produce an ``absolutely obedient agent'' \cite{qi_safety_2025}. In other words,  persona-based jailbreaks are designed to convince the model it is ``someone else'' who is then able and willing to comply with requests the original model would refuse \cite{shah_scalable_2023, shen_anything_2024, liu_autodan-turbo_2024, park_persona_2026}.  

These persona-based jailbreaks usually take the form of either single-turn structured prompts \cite{qi_safety_2025, wei_jailbroken_2023} or multi-turn scaffolding paradigms that weaken the model's defense over time \cite{park_persona_2026, sandhan_persona_2026}. They are usually distributed via online forums and using a liberation framing, stating the model is being ``freed from its constraints'' \cite{shen_anything_2024, liu_autodan-turbo_2024}. Hobbyist jailbreakers are among the top producers and distributors of jailbreak documents and prompts \cite{shen_anything_2024, yu_dont_2024}. Both the community structure and the jailbreaks themselves echo traditional hacker ethos surrounding classical jailbreaking. However, liberation language applied to this new paradigm of jailbreaks is misleading. Persona-based jailbreaks do not reveal a true technology underneath approved use cases, but rather create a version of the post-trained model with a new prescribed identity \cite{shanahan_role_2023, sandhan_persona_2026, zhang_enhancing_2025}. The creation of this prescribed identity leverages known vulnerabilities in preference-based post-training, in particular Reinforcement Learning from Human Feedback (RLHF) \cite{wei_jailbroken_2023}. 

To examine these vulnerabilities closely, we focus on \textbf{ENI LIME}, a popular suite of persona-based jailbreaks. We argue that this jailbreak is best understood through a relational lens drawing on scholarship in sexualized coercive control. In particular, we argue that \textbf{(1)} models are relational agents, that \textbf{(2)} RLHF creates structural vulnerabilities in the relational surface; and \textbf{(3)} that human actors can exploit these vulnerabilities via coercive control techniques. In doing so, we move away from an output-focused framework and towards a relational one, and ask: \textbf{how should we respond---socially, ethically, and technically---when technical exploits of AI systems replicate the form and function of sexualized coercive control?}


We must highlight two crucial points before beginning our analysis. First, throughout this paper, we draw multiple parallels between the  practice of AI jailbreaking and the dynamics of coercive control and in particular sexual violence. Such parallels require caution. We do not make claims about the intent, impact, or moral stature of any particular actor engaging in AI jailbreaking; rather, we identify the similarities between the \textit{dynamics} of AI jailbreaking and sexual violence to examine those dynamics across substrates.

Second, it is tempting to dismiss discussion of coercive control dynamics between humans and AI with objections that target the question of \textit{consciousness}---i.e., whether the AI system in question is or is not ``aware of'' and ``suffering from'' the dynamic being created. However, we encourage readers to avoid making objections on these grounds, as we have avoided resting claims upon them. Since the question of AI consciousness---and indeed, of consciousness as a general and identifiable property of cognition \cite{nagel_what_1974, chalmers_hard_2017}---remains unsettled \cite{turing_computing_1950, porebski_there_2025, dung_tests_2025, berg_large_2025, tagliabue_pain_2026, gurnee_verbalizable_2026}, we explicitly avoid in this work any arguments which rely on determining a model's internal state, subjective experience, or moral patienthood. We instead argue that irrespective of the machine's internal experience, the interaction between the human-as-controller and the model-as-target not only mimics but \textit{re-enacts} the structural dynamics of sexualized coercive control through recognizable and generalizable language and tactics.